% In Catilinam III --- headnote
% in-catilinam-3, 3 December 63 BC

\textit{The third Catilinarian, delivered to the people in
the Forum on the evening of 3 December 63 BC --- the day on
which the conspiracy was at last manifestly caught. The
narrative is the most dramatic in the corpus: through the
Allobrogan envoys, who had been approached by the
conspirators with proposals for a Gallic rising, Cicero had
laid an ambush at the Mulvian Bridge. The praetors Lucius
Flaccus and Gaius Pomptinus, posted in two parts on either
side of the Tiber, intercepted the envoys' party near dawn;
the letters from Lentulus, Cethegus, Statilius, and Gabinius
to the Allobrogan senate and to Catiline were taken with the
seals unbroken. Brought before the Senate that morning, the
conspirators were confronted with their own seals and hands
and confessed in turn. Lentulus, the urban praetor and
self-proclaimed third Cornelius (after Cinna and Sulla)
fated to rule, broke down when shown the seal that bore the
image of his grandfather. Gaius Cethegus, when his cache of
swords and daggers was produced, claimed at first to be a
fancier of fine ironwork, then fell silent. The Senate
decreed custody for the nine principal conspirators;
Lentulus, who had to abdicate his praetorship before he could
be held, did so. Cicero received the first \textit{supplicatio
togato} --- thanksgiving to the gods in the name of a togate
citizen, not a general --- in the recorded history of the
city: ``because I had freed the city from fires, the citizens
from slaughter, Italy from war.''}

\textit{Cicero's narrative occupies more than half the speech.
The closing third develops two themes. The first is the role
of the immortal gods: the lightning-strikes and earthquakes
under Cotta and Torquatus, and the haruspices' instruction to
remake the statue of Jupiter on the Capitoline larger and
turned to the east --- which statue, set up that very morning
by coincidence as the conspirators were being marched in
chains across the Forum to the temple of Concord, faced its
own discovery. ``He, he, Jupiter resisted; he willed the
Capitoline, he these temples, he the entire city, he you
all, to be safe.'' The second is the meditation on what
Cicero has done in comparison with previous civil
dissensions: Sulpicius and Marius, Octavius and Cinna, Lepidus
and Catulus --- all those wished only to alter the
commonwealth, not destroy it; this conspiracy alone willed the
extinction of the city itself, and was conquered ``without
slaughter, without blood, without an army, without battle''
--- by a togate consul. The closing pivot looks ahead, in
hope and in foreboding, to the future of the man who saved
the city: he will live with those whom he has conquered.
The Fourth Catilinarian, two days later, debates what to do
with the prisoners.}
